\documentclass[11pt]{article}
\usepackage[margin=0.72in]{geometry}
\usepackage{amsmath,amssymb,booktabs}
\usepackage[T1]{fontenc}
\usepackage{lmodern}
\usepackage{microtype}
\setlength{\parindent}{0pt}
\setlength{\parskip}{5pt}
\setlength{\abovedisplayskip}{7pt}
\setlength{\belowdisplayskip}{7pt}
\pagestyle{empty}
\newcommand{\area}{\operatorname{area}}
\newcommand{\length}{\operatorname{length}}
\newcommand{\interior}{\operatorname{int}}

\begin{document}
\begin{center}
  {\Large\bfseries Local Voronoi Weight Balancing}
\end{center}
\vspace{-5pt}

\textbf{Setting.}
Let $P\subset\mathbb R^2$ be a compact convex polygon with nonempty
interior, let $C=(c_1,\ldots,c_n)$ consist of distinct centers in
$\interior P$, and let $D\subset P$ be a finite set of dots.
A single step concerns a given center $c_i$.

\medskip
\textbf{1. Cells and neighbors.}
Calculate the bounded Voronoi diagram and the local index set:
\begin{align*}
V_j(C)&=\{x\in P:\|x-c_j\|\leq\|x-c_k\|\text{ for every }k\},
\\
N_i&=\{j\ne i:\length(V_i(C)\cap V_j(C))>0\},
\qquad S_i=\{i\}\cup N_i.
\end{align*}

\textbf{2. Densities and weights.}
For each of these cells $j\in S_i$, let $A_j=\area(V_j(C))$ and let
$M_j$ be its dot count. Assign a dot to its closest center, breaking
ties by the smallest center index. Define
\[
\rho_j=\frac{M_j}{A_j},\qquad W_j=\rho_j A_j=M_j.
\]

\textbf{3. New center location.}
Let $C_i(q)$ be $C$ with $c_i$ replaced by $q$, and define
\[
A_j(q)=\area(V_j(C_i(q))),\qquad
E_i(q)=\frac12\sum_{j\in S_i}\rho_j A_j(q)^2.
\]
Assume the functions $A_j(q)$ are differentiable at $q=c_i$.
For $j\in N_i$, let $\ell_{ij}$ and $m_{ij}$ be the length and
midpoint of the shared edge $V_i(C)\cap V_j(C)$. Set
\[
b_{ij}=\frac{\ell_{ij}}{\|c_i-c_j\|}(m_{ij}-c_i),\qquad
g_i=\nabla E_i(c_i)=\sum_{j\in N_i}(W_i-W_j)b_{ij}.
\]
If $g_i=0$, set $c_i'=c_i$. Otherwise, define the unit direction
and the initial area-change rates by
\[
u_i=-\frac{g_i}{\|g_i\|},\qquad
a_i=\sum_{j\in N_i}b_{ij}\cdot u_i,\qquad
a_j=-b_{ij}\cdot u_i\quad(j\in N_i).
\]
Using the linearized areas $A_j(c_i+s u_i)\approx A_j+s a_j$,
the proposed distance is
\[
\boxed{\displaystyle
s_0=\underset{s\geq0}{\operatorname{arg\,min}}\;
\frac12\sum_{j\in S_i}\rho_j(A_j+s a_j)^2
=\frac{\|g_i\|}{\displaystyle\sum_{j\in S_i}\rho_j a_j^2}.}
\]
Calculate
\[
\boxed{c_i'=c_i+s u_i,\qquad 0\leq s\leq s_0.}
\]
Use the proposed distance $s_0$, stopping earlier if necessary to
keep the center in $\interior P$, avoid reaching another center on
the forward ray $\{c_i+t u_i:t>0\}$, or avoid acquiring a new neighbor.

\newpage
\begin{center}
  {\Large\bfseries Trial and Gap Measures}
\end{center}
\vspace{-5pt}

\textbf{Region and initialization.}
Let $M$ be the existing Manhattan region. Use
\[
P=\operatorname{conv}(M)
\]
and $n=100$ regularly spaced initial centers in $\interior P$.

\textbf{Rounds.}
Run $25$ rounds. In each round, independently draw a uniformly random
permutation of $\{1,\ldots,100\}$ and apply the single-center step in
that order, once for each center.

\textbf{Local gaps.}
For each recorded configuration, calculate the current Voronoi cells,
their neighbor sets $N_i$, and their dot counts $W_i=M_i$. Define
\[
d_{ij}=|W_i-W_j|\qquad(j\in N_i).
\]
For each cell $i$, calculate
\begin{align*}
G_i^{\max}&=\max_{j\in N_i}d_{ij},
\\[4pt]
G_i^{\mathrm{avg}}&=\frac{1}{|N_i|}\sum_{j\in N_i}d_{ij},
\\[4pt]
G_i^{\mathrm{med}}&=\operatorname{median}_{j\in N_i}d_{ij}.
\end{align*}

\textbf{Six measures across cells.}
\begin{center}
\renewcommand{\arraystretch}{1.7}
\begin{tabular}{@{}lccc@{}}
\toprule
 & Maximum gap & Average gap & Median gap \\
\midrule
Maximum
 & $\max_i G_i^{\max}$
 & $\max_i G_i^{\mathrm{avg}}$
 & $\max_i G_i^{\mathrm{med}}$ \\
Median
 & $\operatorname{median}_i G_i^{\max}$
 & $\operatorname{median}_i G_i^{\mathrm{avg}}$
 & $\operatorname{median}_i G_i^{\mathrm{med}}$ \\
\bottomrule
\end{tabular}
\end{center}
For every median, use the middle value for an odd number of entries
and the arithmetic mean of the two middle values for an even number.

\textbf{Recording.}
Record all six measures initially and after every completed round,
$r=0,1,\ldots,25$. Gaps are measured in numbers of dots.

\medskip
\textit{Constraint implementation and a convergence-based stopping rule
remain to be specified.}
\end{document}
